\chapter{Introduction}
\label{ch:introduction}

\section{Background of the Study}
\label{sec:background}

A central question in financial economics and accounting valuation is whether the empirical relationship between published financial statement numbers and equity market prices remains stable across structural corporate reorganizations. At its theoretical core, the valuation of common equity connects secondary market share prices to investors' expectations regarding a firm's future economic earning power, cash-flow generation capability, and net asset backing, conditioned upon the public and private information available in the marketplace \citep{beaver1968,fama1970}. Among competing valuation architectures, the residual income valuation (RIV) framework pioneered by \citet{ohlson1995} and \citet{feltham1995} provides an analytically rigorous and mathematically closed foundation linking financial accounting numbers directly to equity share prices. Anchored in the clean-surplus relation---which dictates that the change in book value of equity equals contemporaneous net income minus net dividends---the Ohlson framework establishes that a firm's equity market value can be expressed as the sum of its currently recognized book value of equity and the discounted present value of expected future residual income (abnormal earnings). Under linear information dynamics, this theoretical formulation demonstrates that market prices can be represented as a linear function of contemporaneous earnings per share (EPS) and book value per share (BVPS). In this framework, BVPS serves as an accounting proxy for the cumulative historical net assets contributed by shareholders and retained from past earnings, whereas EPS captures current operational profitability and serves as the baseline indicator of the firm's forward-looking earning capacity \citep{collins1997}.

Empirical studies investigating the association between equity prices and published accounting fundamentals---collectively recognized as the accounting value-relevance literature---have provided extensive documentation across developed, emerging, and frontier equity markets \citep{almumani2014,malhotra2013,sharma2011}. In frontier and developing market environments, however, the informational capacity of equity markets is frequently challenged by structural friction, regulatory evolution, asymmetric information, illiquidity, and corporate restructuring. The Nepalese capital market presents a distinct and compelling setting for examining the empirical validity and temporal stability of accounting value relevance. The organized secondary securities market in Nepal originated with the establishment of the Securities Exchange Centre in 1976 (2033 BS), which operated under the Securities Exchange Act, 1983. In 1993, a major financial-sector reform initiative restructured the Centre into the Nepal Stock Exchange (NEPSE) Limited as the nation's sole organized stock exchange, concurrently establishing the Securities Board of Nepal (SEBON) as the apex statutory securities market regulator. Floor trading on NEPSE officially commenced on 13 January 1994 (29 Poush 2050 BS) \citep{nepse2017}. Since its inception, NEPSE has been heavily dominated by the banking and financial institutions (BFIs) sector, with Class 'A' commercial banks constituting a major component of aggregate market capitalization, total annual trading turnover, and public equity ownership \citep{pradhan2016,timalsina2026}.

Nepal's commercial banking sector has experienced profound structural, regulatory, and institutional transformations over the past two decades. Following the financial liberalization policies initiated in the late 1980s, the number of licensed commercial banks expanded rapidly, reaching an apex of 32 licensed commercial banks alongside dozens of development banks and finance companies by fiscal year 2011/12 according to Nepal Rastra Bank institutional records \citep{adhikari2023,nrb2011}. Recognizing that severe institutional proliferation was outstripping supervisory capacity, inducing cutthroat lending competition, and creating systemic vulnerability, Nepal Rastra Bank (NRB)---the central bank and banking regulator---embarked on an aggressive, multi-phase consolidation drive. In 2011 (2068 BS), NRB introduced the \textit{Bank and Financial Institutions Merger and Acquisition Bylaws}, providing the initial regulatory architecture for voluntary corporate combinations \citep{nrb2011}. Subsequently, in its Monetary Policy for FY 2015/16 (2072/73 BS), NRB executed a historic fourfold capital expansion mandate, requiring all Class 'A' commercial banks to increase their minimum paid-up equity capital from NPR 2.0 billion to NPR 8.0 billion within a two-year regulatory window \citep{nrb2015}. Although banks initially satisfied this capital mandate primarily through bonus share capitalization and rights offerings, the persistent pressure for systemic consolidation culminated in the major banking consolidation wave of 2022--2023 (2079 BS). Supported by regulatory incentives and moral suasion, prominent commercial banks executed large-scale horizontal amalgamations and statutory acquisitions. By mid-October 2023 (Asoj 2080 BS), successive corporate consolidations reduced the commercial bank population to exactly 20 licensed institutions \citep{nrb2023list}.

A central theoretical and empirical issue in financial accounting is whether the valuation relationship between accounting fundamentals and equity market prices remains stable across major corporate reorganizations and structural transitions. As demonstrated in international M\&A literature, large-scale corporate mergers fundamentally restructure the corporate reporting entity: they alter the financial scale of the institution, introduce substantial post-merger integration costs, require accounting harmonization across legacy loan and deposit portfolios, generate non-recurring organizational adjustments, and expand the aggregate share capital base \citep{kwon2020}. Under such restructuring, reported accounting earnings can experience temporary impairment in transitoriness, volatility, and noise, potentially leading equity investors to discount earnings and rely more heavily on balance-sheet equity (book value) as a valuation anchor \citep{collins1997}. In Nepal, four of the country's most prominent commercial banking series---Global IME Bank Limited (GBIME), Himalayan Bank Limited (HBL), Kumari Bank Limited (KBL), and Nepal Investment Mega Bank Limited (NIMB)---commenced joint operations following amalgamations between January and February 2023 (Poush--Falgun 2079 BS). Concurrently, established peer institution NMB Bank Limited (NMB) maintained operational continuity without participating in the early-2023 consolidation wave, thereby providing a contemporaneous benchmark series. Although NMB had previously acquired Kanchan Development Bank Limited in August 2020 (adding approximately NPR 0.56 billion beyond its 13\% bonus share distribution to its share capital in FY 2077/78), it did not undergo horizontal amalgamation during the 2022--2023 consolidation wave, remaining structurally stable throughout the focal restructuring divide.

Despite the profound economic and corporate restructuring induced by the 2023 consolidation wave, empirical research on Nepalese commercial banking valuation reviewed in Chapter II has predominantly focused on static, sample-average regressions spanning multi-year historical periods without explicitly evaluating whether the fundamental accounting-price relationships shifted across the consolidation divide \citep{bhattarai2014,chaudhari2025,chhetri2023,gurungdahal2025,joshi2023,pant2024,timalsina2026}. Furthermore, conventional literature has routinely suffered from methodological shortcomings, notably the circular inclusion of price-derived valuation ratios such as Price-to-Earnings (P/E) as explanatory regressors predicting market prices, and the reliance on fiscal-year-end (Ashadh-end) pricing dates that precede the public disclosure of audited financial accounts. This dissertation directly addresses these theoretical, institutional, and empirical gaps by evaluating the value relevance of accounting information (EPS and BVPS) before and after Nepal's 2023 banking consolidation wave across a longitudinal panel of five commercial bank series over a six-year period (FY 2076/77 to FY 2081/82), using a standardized post-reporting-deadline pricing window grounded in Section 59(2) of the Banks and Financial Institutions Act, 2017 \citep{bafia2017}.

Previewing the empirical findings, the baseline two-way fixed-effects model yields positive but statistically insignificant slopes for both accounting fundamentals: earnings per share ($\hat{\beta}_E = 2.095$, HC3 $t = 0.51$, $p = 0.616$, with $16$ residual degrees of freedom) and book value per share ($\hat{\beta}_B = 0.219$, HC3 $t = 0.12$, $p = 0.908$). In the interaction specification (Model~2, with $14$ residual degrees of freedom), the pre-consolidation earnings slope is economically substantial but not statistically significant ($\hat{\beta}_E = 5.577$, $p = 0.140$), while the pre-consolidation book value slope is negligible ($\hat{\beta}_B = 0.100$, $p = 0.940$). The estimated post-consolidation shift parameters---$\hat{\delta}_E = -3.124$ ($p = 0.448$) and $\hat{\delta}_B = -1.758$ ($p = 0.175$)---imply compressed post-consolidation slopes of $\hat{\beta}_E + \hat{\delta}_E = 2.453$ for earnings and $\hat{\beta}_B + \hat{\delta}_B = -1.658$ for book value, but neither individual shift attains significance. The joint Wald test of structural change indicates parameter shift across regimes at the 10\% level ($F(2, 14) = 2.954, p = 0.085$), with structural instability confirmed at the 5\% level when the FY 2078/79 entity-timing mismatch is excluded ($F(2, 10) = 5.818, p = 0.021$).

\section{Problem Statement}
\label{sec:problem_statement}

The valuation behavior of commercial bank equities in emerging capital markets is governed by a complex interplay between institutional regulation, financial disclosure credibility, corporate scale, and investor sentiment. In Nepal, empirical investigations examining the association between accounting indicators and equity share prices have documented extensive but divergent findings. Prior studies have evaluated diverse financial measures, including earnings per share, book value per share, dividend per share, return on assets, return on equity, and price-earnings ratios, yet their reported empirical conclusions diverge sharply regarding coefficient magnitudes, directional signs, and statistical significance.

For instance, regarding operating profitability, \citet{chhetri2023} examined 13 commercial banks over 2012--2022 using pooled ordinary least squares (OLS) and reported a statistically significant positive association between return on assets (ROA) and market share price (MPS). Conversely, \citet{chaudhari2025} investigated ten commercial banks over 2014/15--2023/24 and documented a statistically significant \textit{negative} relationship between ROA and share prices, while \citet{gyawali2022} found a positive but statistically insignificant ROA coefficient across ten banks observed between 2017 and 2021. For book value per share, \citet{joshi2023} reported a strong, positive, statistically significant association with market price among seven commercial banks over 2012/13--2021/22, whereas \citet{pant2024} documented an insignificant negative coefficient for BVPS across seven commercial banks observed through FY 2023/24. Similarly, while \citet{bhandari2026} identified EPS as the predominant determinant of equity valuation across listed commercial banks, \citet{chhetri2023} and \citet{giri2024} found bivariate or multivariate EPS coefficients to be statistically negligible or secondary to dividend and net worth measures.

These empirical discrepancies stem fundamentally from differences in sample composition, observation periods, econometric specifications, and variable measurement timing. However, beyond empirical divergence across static studies, the existing literature reveals four critical conceptual and methodological problems that motivate the present dissertation:

\begin{enumerate}[label=\textbf{\arabic*.}]
\item \textbf{The Absence of Consolidation Regime Shift Analysis:}
Published Nepalese studies reviewed in this dissertation predominantly employ pooled OLS or static panel regression designs that estimate unconditional, full-sample average coefficients across their entire observation spans \citep{bhattarai2014,chaudhari2025,gurungdahal2025,timalsina2026}. These static models implicitly assume parameter stationarity---that the valuation multiple attached to a rupee of earnings or book value remains identical throughout the sample period. However, between January and February 2023, four prominent commercial banking mergers in recent Nepalese history commenced joint operations: Global IME Bank absorbed Bank of Kathmandu (9 January 2023); Himalayan Bank acquired Civil Bank (24 February 2023); Kumari Bank absorbed NCC Bank (1 January 2023); and Nepal Investment Bank amalgamated with Mega Bank Nepal to form NIMB (11 January 2023) \citep{gbime2022agm,hbl2023report,kumari2022sgm,nimb2023report}. These amalgamations fundamentally transformed the capital scale, loan portfolios, branch networks, and risk profiles of the surviving entities. International research demonstrates that structural corporate combinations alter the financial reporting environment, temporarily increasing earnings volatility while shifting investor attention toward balance-sheet net worth \citep{kwon2020}. No empirical study identified in the literature reviewed in Chapter II has formally evaluated whether the valuation coefficients of EPS and BVPS shifted between the pre-consolidation (FY 2076/77--2078/79) and post-consolidation (FY 2079/80--2081/82) periods using an explicit regime-interaction econometric framework.

\item \textbf{Econometric Endogeneity from Price-Derived Explanatory Variables:}
A common practice across studies reviewed in Chapter II is the inclusion of the price-earnings (P/E) ratio alongside fundamental accounting variables in regression models where Market Price per Share (MPS) serves as the dependent variable \citep{bhattarai2014,chhetri2023,giri2024,gurungdahal2025,joshi2023,pant2024}. Because the P/E ratio is algebraically defined as $\text{P/E} = \text{MPS} / \text{EPS}$, specifying a model of the form $\text{MPS}_{it} = \beta_0 + \beta_1 \text{EPS}_{it} + \beta_2 \text{P/E}_{it} + \dots$ embeds the dependent variable directly within the numerator of one of its own explanatory regressors. This introduces mechanical endogeneity, circular causality, and spurious correlation that severely distort the estimated coefficients and invalidate standard hypothesis testing. A methodologically defensible accounting valuation study must strictly exclude price-derived multiples and focus exclusively on fundamental accounting variables that are determined independently of secondary market trading prices.

\item \textbf{Temporal Misalignment and Look-Ahead Bias in Price Measurement:}
A critical limitation in prior Nepalese valuation research concerns the timing of share price measurement. Across the empirical studies reviewed in Chapter II, published literature predominantly uses fiscal-year-end (Ashadh-end / mid-July) closing market prices or annual volume-weighted average prices to evaluate the relevance of financial statement numbers \citep{bhattarai2014,pradhan2016,silwal2019}. Under Nepalese corporate and regulatory realities, commercial bank fiscal years terminate at Ashadh-end (mid-July), but audited annual financial statements are not prepared, audited by statutory auditors, approved by NRB, and published for months thereafter. Pairing Ashadh-end market prices with audited accounting figures that were not publicly known on that date creates severe look-ahead bias, violating the temporal causality required for information content studies \citep{ball1968,beaver1968,bernard1989}. Conversely, attempting to reconstruct individual, ad hoc corporate disclosure dates for 30 bank-year observations is fraught with measurement noise, selective informal leakage, and archival opacity. To resolve this temporal conundrum, an objective, legally grounded measurement rule is required. Section 59(2) of the Banks and Financial Institutions Act, 2017 (BAFIA) legally mandates that Class 'A' commercial banks publish their audited annual balance sheets and profit and loss accounts within nine months into the subsequent fiscal year---a statutory deadline operationalized at the end of Chaitra (mid-April) of the subsequent fiscal year \citep{bafia2017}. Observing market prices at a standardized window immediately following this statutory threshold ensures that audited financial information was legally required to be public, avoiding both look-ahead bias and timestamp reconstruction ambiguity.

\item \textbf{Neglect of Degrees of Freedom and Small-Cluster Econometric Fragility:}
Prior single-sector panel studies in Nepal frequently implement econometric estimation without disclosing or addressing the severe statistical constraints imposed by small cross-sectional dimensions. In fixed-effects panel models with interactive specifications, parameter proliferation rapidly consumes residual degrees of freedom. On a clean sample of $N = 28$ observations, estimating individual bank fixed effects, common fiscal-year fixed effects, and four slope coefficients leaves exactly $28 - 14 = 14$ residual degrees of freedom. Furthermore, standard cluster-robust sandwich standard errors perform poorly and exhibit substantial finite-sample downward bias when the number of clusters is small ($G = 5$) \citep{cameron2008clusters}. Prior empirical studies in Nepal reviewed in Chapter II have largely left these inferential hazards unaddressed, risking false discoveries and overstated statistical precision. Rigorous modeling requires explicit parameter accounting and the deployment of MacKinnon-White HC3 heteroskedasticity-robust inference, accompanied by transparent non-causal interpretation.
\end{enumerate}

To address these empirical, methodological, and econometric problems, this study investigates three central research questions:
\begin{enumerate}[label=\textbf{\arabic*.}]
\item \textbf{Research Question 1:} Are fundamental accounting measures---specifically earnings per share (EPS) and book value per share (BVPS)---statistically associated with post-reporting-deadline market share prices ($MPS^{RD}$) among the selected Nepalese commercial banks?
\item \textbf{Research Question 2:} Did the estimated valuation associations of EPS and BVPS with market share prices differ significantly between the pre-consolidation period (FY 2076/77--2078/79) and the post-consolidation period (FY 2079/80--2081/82)?
\item \textbf{Research Question 3:} Are the baseline and consolidation-interaction findings robust to alternative operating profitability specifications, inclusion of eligible dividend disclosures, adjustment for bonus share capitalization, retrospective earnings restatements, and the econometric exclusion of structural transition periods?
\end{enumerate}

\section{Objectives of the Study}
\label{sec:objectives}

The overarching general objective of this dissertation is to examine the accounting value relevance of financial information and evaluate whether the empirical relationships between accounting fundamentals and equity market prices differed before and after the 2023 commercial banking consolidation wave across five purposively selected commercial bank series in Nepal.

To achieve this general objective, the study establishes the following specific objectives:
\begin{enumerate}[label=\textbf{\arabic*.}]
\item \textbf{To estimate the baseline value relevance} of earnings per share (EPS) and book value per share (BVPS) in relation to secondary market share prices measured at a standardized post-reporting-deadline window ($MPS^{RD}$) within a two-way fixed-effects panel regression framework.
\item \textbf{To test whether the empirical valuation slopes} of EPS ($\beta_E$) and BVPS ($\beta_B$) differed between the pre-consolidation period (FY 2076/77--2078/79) and the post-consolidation period (FY 2079/80--2081/82) by evaluating interaction parameters ($\delta_E, \delta_B$) and a joint Wald coefficient test.
\item \textbf{To evaluate the econometric robustness and stability} of the value-relevance estimates across alternative specifications, including:
\begin{itemize}
\item Substituting return on assets (ROA) for EPS as an asset-scaled profitability measure;
\item Incorporating eligible cash dividend per share (DPS) declarations into the valuation equation;
\item Sensitivity analysis excluding the FY 2079/80 merger transition year to verify results are not driven by immediate consolidation shock;
\item Sensitivity analysis excluding FY 2078/79 to address entity-timing mismatch between standalone fundamentals and consolidated trading prices;
\item Re-estimation using bonus-adjusted per-share metrics to account for share capital dilutions; and
\item Re-estimation using retrospective audited basic EPS restatements from subsequent annual reports.
\end{itemize}
\end{enumerate}

\section{Hypotheses of the Study}
\label{sec:hypotheses}

The research hypotheses formulated in this dissertation are directly anchored in the residual income valuation theory of \citet{ohlson1995}, the post-merger value-relevance framework of \citet{kwon2020}, and empirical findings from emerging capital markets. Under Ohlson's (1995) theoretical model, both contemporaneous earnings and book equity are expected to correlate positively with market prices because they reflect, respectively, current operating efficiency and accumulated net asset value. However, following the methodological conventions of empirical accounting, all hypotheses are tested two-tailed at the 5\% level of significance ($\alpha = 0.05$) without imposing directional signs, thereby allowing the empirical data to support or fail to support the theoretical predictions.

The hypotheses are grouped into two categories corresponding to the baseline value-relevance model (Model 1) and the consolidation-interaction model (Model 2):

\subsection{Category A: Baseline Value Relevance Hypotheses (Model 1)}
\begin{itemize}
\item \textbf{Hypothesis 1 ($H_1$):}
\begin{itemize}
\item Null hypothesis ($H_{0,1}$): Earnings per share (EPS) is not significantly associated with secondary market share price at the reporting-deadline window in the baseline two-way fixed-effects model.
\begin{equation}
H_{0,1}: \beta_E = 0
\end{equation}
\item Alternative hypothesis ($H_{a,1}$): Earnings per share (EPS) is significantly associated with secondary market share price at the reporting-deadline window in the baseline two-way fixed-effects model.
\begin{equation}
H_{a,1}: \beta_E \ne 0
\end{equation}
\end{itemize}

\item \textbf{Hypothesis 2 ($H_2$):}
\begin{itemize}
\item Null hypothesis ($H_{0,2}$): Book value per share (BVPS) is not significantly associated with secondary market share price at the reporting-deadline window in the baseline two-way fixed-effects model.
\begin{equation}
H_{0,2}: \beta_B = 0
\end{equation}
\item Alternative hypothesis ($H_{a,2}$): Book value per share (BVPS) is significantly associated with secondary market share price at the reporting-deadline window in the baseline two-way fixed-effects model.
\begin{equation}
H_{a,2}: \beta_B \ne 0
\end{equation}
\end{itemize}
\end{itemize}

\subsection{Category B: Consolidation-Period Regime Shift Hypotheses (Model 2)}
In corporate reorganizations, international research suggests that post-merger integration costs, portfolio revaluation, and organizational disruption may temporarily increase earnings noise, potentially leading investors to place relatively greater valuation weight on balance-sheet book value \citep{kwon2020}. However, whether this shift manifested in Nepal's emerging banking sector remains an open empirical question. Accordingly, the interaction hypotheses test for parameter differences across periods:
\begin{itemize}
\item \textbf{Hypothesis 3 ($H_3$):}
\begin{itemize}
\item Null hypothesis ($H_{0,3}$): The slope coefficient linking earnings per share to market share price does not differ significantly between the pre-consolidation and post-consolidation periods.
\begin{equation}
H_{0,3}: \delta_E = 0
\end{equation}
\item Alternative hypothesis ($H_{a,3}$): The slope coefficient linking earnings per share to market share price differs significantly between the pre-consolidation and post-consolidation periods.
\begin{equation}
H_{a,3}: \delta_E \ne 0
\end{equation}
\end{itemize}

\item \textbf{Hypothesis 4 ($H_4$):}
\begin{itemize}
\item Null hypothesis ($H_{0,4}$): The slope coefficient linking book value per share to market share price does not differ significantly between the pre-consolidation and post-consolidation periods.
\begin{equation}
H_{0,4}: \delta_B = 0
\end{equation}
\item Alternative hypothesis ($H_{a,4}$): The slope coefficient linking book value per share to market share price differs significantly between the pre-consolidation and post-consolidation periods.
\begin{equation}
H_{a,4}: \delta_B \ne 0
\end{equation}
\end{itemize}
\end{itemize}

The overarching econometric evaluation of structural change across the banking consolidation wave is conducted via the joint null and alternative hypotheses:
\begin{equation}
H_{0,\text{joint}}: \delta_E = 0 \text{ and } \delta_B = 0 \quad \text{versus} \quad H_{a,\text{joint}}: \delta_E \ne 0 \text{ or } \delta_B \ne 0
\end{equation}
Rejecting this joint null hypothesis provides statistical evidence that the overall accounting-price valuation relationship differed across the consolidation period. In accordance with rigorous observational econometrics, non-rejection or rejection of these parameters is interpreted strictly as observed differences in statistical associations across calendar regimes, rather than as causal impacts of bank mergers.

\section{Rationale and Significance of the Study}
\label{sec:rationale}

This dissertation is justified on theoretical, methodological, empirical, and practical grounds. As Nepal's commercial banking sector emerges from the structural consolidation wave of 2022--2023, understanding how capital markets perceive and price accounting disclosures is critical for financial transparency, capital allocation, and regulatory design. The specific rationale and significance of this research include:

\begin{enumerate}[label=\textbf{\arabic*.}]
\item \textbf{Theoretical and Conceptual Contribution:}
While the residual income valuation framework of \citet{ohlson1995} has been extensively tested in developed equity markets, its applicability in frontier environments undergoing rapid structural transformation remains underexplored. By synthesizing Ohlson's (1995) valuation model with the post-merger value-relevance framework of \citet{kwon2020}, this study provides a coherent theoretical structure for assessing how corporate amalgamations alter the relative informational roles of income-statement flows (earnings) versus balance-sheet stocks (book value). It grounds the empirical valuation of Nepalese banks within formal accounting valuation theory.

\item \textbf{Methodological Rigor and Standardization of Observation Timing:}
A major contribution of this study is the development and implementation of an institutionally grounded, reproducible price-measurement protocol ($MPS^{RD}$). By operationalizing the statutory nine-month reporting deadline under Section 59(2) of BAFIA at the end of Chaitra (mid-April), the study addresses the temporal flaw of prior Nepalese studies that paired Ashadh-end market prices with financial reports published up to nine months later. In observed corporate practice across the sampled banks, audited financial highlights and dividend declarations are typically published between Ashwin and Poush (October to January), with formal AGM proceedings concluded prior to the end-of-Chaitra reporting deadline operationalized in this study. Observing market prices at the first trading session following the Chaitra deadline ensures that audited financial accounts were legally required to be in the public domain under BAFIA Section 59(2) prior to price observation, while avoiding the measurement error and selective leakage inherent in attempting to reconstruct unverified corporate press-release timestamps across historical periods. Furthermore, the explicit formulation of a 5-day freshness rule---under which trading-suspended observations are treated as missing rather than imputed from stale dates---maintains empirical transparency and eliminates stale price contamination.

\item \textbf{Econometric Integrity through Exclusion of Mechanical Multiples:}
This dissertation explicitly corrects the pervasive econometric specification error in Nepalese literature: the inclusion of P/E ratios alongside EPS in models predicting share prices. By excluding price-derived multiples and demonstrating the severe collinearity between EPS and ROA ($r = 0.975$, yielding variance inflation factors between approximately 20 and 30 depending on model specification), this study avoids the mechanical simultaneity bias inherent in price-deflated multiple regressions. It demonstrates why ROA must be evaluated as an alternative substitute rather than an additive regressor, preventing spurious statistical inferences.

\item \textbf{Practical Significance for Capital Market Stakeholders:}
The findings of this study provide actionable insights for diverse capital market participants:
\begin{itemize}
\item \textbf{Individual and Institutional Investors:} Informs investors on NEPSE regarding the comparative valuation relevance of earnings versus book values across structural consolidation periods, highlighting the valuation noise associated with transitional merger accounting.
\item \textbf{Commercial Bank Management and Boards of Directors:} Informs bank executives regarding the capital-market consequences of earnings volatility, dividend payout stability, and balance-sheet transparency following corporate amalgamations.
\item \textbf{Nepal Rastra Bank (NRB):} Offers the central bank empirical observations on how secondary equity markets evaluate capital and earnings metrics surrounding policy-driven bank consolidations, assisting in the design of future prudential capital and reporting guidelines.
\item \textbf{Securities Board of Nepal (SEBON) and NEPSE:} Highlights the critical importance of timely corporate disclosures, centralized announcement timestamp databases, and the market-distorting effects of prolonged merger-related trading suspensions.
\item \textbf{Academic Researchers:} Provides a documented, reproducible dataset and analytical codebase, rigorous two-way fixed-effects methodology, and robust finite-sample econometric protocols for future researchers investigating financial-sector dynamics in South Asia.
\end{itemize}
\end{enumerate}

\section{Limitations of the Study}
\label{sec:limitations}

To maintain academic integrity and assist readers in interpreting the findings within their proper context, the methodological, institutional, and empirical boundaries of this research are explicitly acknowledged. In accordance with the \textit{Master's Degree Dissertation Guidelines (2019)} of the Faculty of Management, a clear distinction is maintained between the scope of the study (which defines its deliberate boundaries) and the limitations that arise within that scope:

\begin{enumerate}[label=\textbf{\arabic*.}]
\item \textbf{Small Purposive Sample and Sectoral Generalizability:}
The study's scope is confined to five purposively selected Class 'A' commercial bank series observed across six fiscal years (FY 2076/77 to FY 2081/82). While these five institutions represent the central corporate combinations of the 2023 consolidation wave alongside an unmerged benchmark (in the 2023 wave), the sample cannot be generalized to the entire commercial banking industry, development banks, finance companies, or non-financial listed sectors. Findings apply directly to the selected institutions and their specific consolidation contexts.

\item \textbf{Realized Sample Attrition Due to Trading Suspensions:}
Although primary audited annual financial statements were complete for all 30 target bank-year observations, prolonged merger-related trading suspensions rendered two pre-consolidation market price observations unverifiable within the 5-day freshness window: Himalayan Bank FY 2077/78 (merger-related trading suspension until 6 June 2022; 54 days late) and Nepal Investment Bank FY 2077/78 (merger-related trading suspension until 30 May 2022; 47 days late). Consequently, following strict research protocol, these two observations are coded as missing rather than imputed from stale pre-suspension prices, reducing the clean estimation sample to $N = 28$ observations (13 pre-consolidation, 15 post-consolidation).

\item \textbf{Econometric Fragility and Finite-Sample Constraints:}
Estimating two-way fixed effects (bank fixed effects $\mu_i$ and fiscal-year fixed effects $\lambda_t$) in conjunction with four slope parameters on a sample of $N = 28$ observations leaves exactly $28 - 14 = 14$ residual degrees of freedom in Model 2. Furthermore, standard cluster-robust standard errors fail to perform reliably when the number of clusters is very small ($G = 5$) \citep{cameron2008clusters}. Although the study mitigates heteroskedasticity by deploying MacKinnon-White HC3 robust standard errors, all statistical tests must be interpreted as exploratory indicators rather than definitive asymptotic proofs.

\item \textbf{Calendar-Period Operationalization vs. Causal Identification:}
The pre-consolidation (FY 2076/77--2078/79) and post-consolidation (FY 2079/80--2081/82) regimes are defined by calendar time. Consequently, the estimated interaction parameters ($\delta_E, \delta_B$) capture overall shifts in accounting-price associations across these temporal periods. These shifts reflect not only the internal corporate effects of banking consolidation, but also contemporaneous macroeconomic shifts, interest rate fluctuations, NEPSE broad-index cycles, and post-pandemic economic adjustments. The research design does not identify the isolated causal impact of mergers.

\item \textbf{Entity-Timing Mismatch in FY 2078/79:}
Because the four merging banks commenced joint operations in January--February 2023 (mid-FY 2079/80), their FY 2078/79 audited financial accounts represent pre-merger standalone operations, whereas their post-reporting-deadline market prices observed in mid-April 2023 reflect post-merger consolidated trading. Although this transition mismatch is retained in the baseline models to preserve sample continuity, it represents an unavoidable institutional friction that is explicitly evaluated via sensitivity exclusions.

\item \textbf{Institutional Reporting Compliance and Information Leakage:}
The standardized reporting window relies on the legal mandate under BAFIA Section 59(2) requiring publication within nine months into the subsequent fiscal year. While in the disclosures examined for this study, commercial bank boards typically released audited highlights and dividend proposals between Ashwin and Poush (months prior to the Chaitra statutory deadline), this pattern reflects observed corporate practice rather than audited proof of disclosure timing for every individual observation; any unrecorded informal information leakage prior to public board announcements remains an unobservable institutional feature of frontier equity markets.
\end{enumerate}

\section{Organization of the Study}
\label{sec:organization}

This dissertation is organized into five chapters, structured in accordance with the \textit{Master's Degree Dissertation Guidelines (2019)} of the Faculty of Management, Tribhuvan University:

\begin{itemize}
\item \textbf{Chapter I: Introduction:} Provides the contextual background of the study, presents the problem statement, outlines the research questions, defines the general and specific objectives, states the research hypotheses, justifies the rationale and significance of the study, acknowledges the research limitations, and outlines the overall chapter structure.
\item \textbf{Chapter II: Literature Review:} Reviews the theoretical underpinnings of equity valuation, including the Ohlson (1995) residual income valuation framework, post-merger value-relevance theory, signaling theory, agency theory, and efficient market dynamics. Systematically reviews international and Nepalese empirical studies on banking valuation, examines the evolution of merger regulations in Nepal, and synthesizes the identified literature gaps that this study addresses.
\item \textbf{Chapter III: Research Methodology:} Explains the quantitative, correlational, and comparative panel research design. Details the target population, purposive sampling criteria, and profile of the five sampled commercial bank series. Specifies the secondary data sources, operational variable definitions ($MPS^{RD}$, EPS, BVPS, ROA, DPS, $\mathit{Post}_t$), and variable measurement rules. Formulates the baseline and consolidation-interaction two-way fixed-effects models, details the MacKinnon-White HC3 econometric estimation procedure, outlines the diagnostic tests, and specifies the eight-specification estimation and robustness suite (R1--R2 baselines plus six robustness variants R3--R8).
\item \textbf{Chapter IV: Results and Discussion:} Presents the empirical findings derived from the analysis pipeline. Details descriptive statistics across full, pre-, and post-consolidation periods; presents the Pearson correlation matrix and screens for multicollinearity; reports parameter estimates and hypothesis test results for Model 1 (baseline) and Model 2 (interaction); presents the joint Wald $F$-test of structural shift; reports the results of the eight-specification estimation and robustness suite; presents econometric diagnostic tests (VIF, Breusch-Pagan, Cook's distance, normality); and provides a structured discussion contextualizing the empirical findings with prior Nepalese and international literature.
\item \textbf{Chapter V: Summary and Conclusion:} Summarizes the entire study, synthesizing the core empirical findings regarding earnings and book value relevance across the consolidation wave. Draws overarching conclusions based on the statistical evidence, outlines actionable policy and managerial implications for investors, bank executives, and regulatory authorities (NRB, SEBON), and offers recommendations for future academic research.
\end{itemize}

Following the main body of the dissertation, a complete \textbf{References} section compiled in APA 6th edition format provides full bibliographic details for all cited literature, statutes, and reports. A comprehensive \textbf{Appendices} section contains supplementary tables (price verification records, specification comparisons, full regression dataset) and appendix diagnostic figures.

